\documentclass[11pt,a4paper]{article}
\usepackage[utf8]{inputenc}
\usepackage[margin=1in]{geometry}
\usepackage{amsmath}
\usepackage{amssymb}
\usepackage{booktabs}
\usepackage{hyperref}
\usepackage{enumitem}
\usepackage{tocloft}
\usepackage{fancyhdr}

% Headers and footers
\pagestyle{fancy}
\fancyhf{}
\rhead{\thepage}
\lhead{\leftmark}
\renewcommand{\headrulewidth}{0.4pt}

% Title formatting
\title{\textbf{Midas: Credit Portfolio Model With Climate Overlay}\\
\large Model for Integrated Default Analysis with Scenarios\\[0.5em]
\normalsize A factor-based framework with deterministic climate/scenario drivers\\and systematic credit factors with market-calibrated correlation structure}
\author{}
\date{\today}

\begin{document}

\maketitle

\tableofcontents
\newpage

\section*{About Midas}

The name \textbf{Midas} stands for \textbf{M}odel for \textbf{I}ntegrated \textbf{D}efault \textbf{A}nalysis with \textbf{S}cenarios. Midas transforms climate scenario data into actionable credit risk insights through a transparent, governance-friendly factor model that preserves familiar portfolio modeling behavior while incorporating climate drivers.

\section{Purpose and scope}

\subsection{What we are trying to achieve}

We want a credit portfolio model that produces a distribution of defaults and losses under a chosen deterministic climate scenario, while preserving familiar factor-model behaviour (correlated outcomes, portfolio tails) using a structured, calibratable correlation layer.

The model should:

\begin{enumerate}
    \item Behave like a standard credit factor model (familiar to risk practitioners and governance).
    \item Translate climate scenario variables into systematic credit stress in a transparent, auditable way.
    \item Support scenario analysis and stress testing rather than claim to ``predict'' climate-driven defaults from first principles.
    \item Be implementable even when we lack clean historical observables to statistically estimate climate-credit parameters.
\end{enumerate}

\subsection{Key simplification we adopt}

We keep one main factor object that drives correlated credit outcomes:
\begin{itemize}
    \item \textbf{F}: a sector--region factor matrix, representing systematic ``credit conditions'' for each sector and region cell.
\end{itemize}

Climate enters through:
\begin{itemize}
    \item \textbf{X}: a deterministic vector of climate/scenario variables (e.g., carbon price path, energy prices, hazard intensities), which drives $F$ via a mapping $F = m(X) + u$.
\end{itemize}

So the model is:
\begin{itemize}
    \item Credit outcomes depend on $F$
    \item $F = m(X) + u$ where $m(X)$ is deterministic scenario mean, $u$ is stochastic residual
    \item Idiosyncratic noise provides name-level randomness.
\end{itemize}

\section{Conceptual overview}

\subsection{Three layers of the system}

\begin{enumerate}
    \item \textbf{Scenario layer (deterministic)}: A climate scenario gives a set of driver values $X$ (e.g., carbon price change, energy price shock, physical hazard indices). These are \textbf{not random} in the model.

    \item \textbf{Systematic credit layer (stochastic)}: The scenario shifts the mean systematic credit conditions for each sector $\times$ region cell. Around that mean, there is a stochastic residual credit environment that creates correlation and tail risk.

    \item \textbf{Obligor layer (stochastic)}: Each obligor loads on the sector $\times$ region systematic credit conditions via weights $\beta_i$ and has an idiosyncratic shock. A single ``asset-correlation'' knob $\rho$ controls how strongly obligors co-move with systematic conditions overall.
\end{enumerate}

\subsection{Why a sector-region cell structure?}

The cell-level factor matrix $F_{s,r}$ provides:
\begin{itemize}
    \item \textbf{Natural interpretation}: $F_{\text{coal,europe}}$ directly represents ``coal sector in Europe''
    \item \textbf{Flexible exposures}: Obligors can load on multiple cells (diversified portfolios)
    \item \textbf{Climate overlay}: Deterministic scenario $X$ drives systematic mean $m_{s,r}(X)$
    \item \textbf{Residual correlation}: Stochastic $u$ captures non-climate systematic uncertainty via market-calibrated covariance $\Sigma_u$
\end{itemize}

\section{Underlying theory: credit portfolio factor models}

\subsection{The classic structure}

Most portfolio credit models start by constructing a latent creditworthiness variable for each obligor $i$:
\begin{equation}
Z_i = \text{systematic part} + \text{idiosyncratic part}
\end{equation}

The systematic part is shared across obligors (causing correlated outcomes). The idiosyncratic part is independent.

A common, governance-friendly form is:
\begin{equation}
Z_i = \sqrt{\rho} \cdot (\text{systematic score}) + \sqrt{1-\rho} \cdot \varepsilon_i
\end{equation}

where:
\begin{itemize}
    \item $\rho \in (0,1)$: fraction of variance that is systematic (``asset correlation'')
    \item $\varepsilon_i$: idiosyncratic noise (standard normal)
\end{itemize}

A default event is triggered by a threshold:
\begin{equation}
\text{Default if } Z_i < \tau_i
\end{equation}

The threshold $\tau_i$ is set to reproduce the obligor's unconditional probability of default (PD), typically via:
\begin{equation}
\tau_i = \Phi^{-1}(\text{PD}_i)
\end{equation}

where $\Phi$ is the standard normal CDF.

\subsection{Why this works}

\begin{itemize}
    \item The model reproduces marginal PDs by construction.
    \item It generates default correlation through shared systematic component.
    \item It can be simulated efficiently (Monte Carlo) and decomposed by risk drivers.
    \item The single parameter $\rho$ controls tail thickness and diversification benefits.
\end{itemize}

\section{Our model specification}

\subsection{Index sets and dimensions}

\begin{itemize}
    \item \textbf{Sectors}: $s = 1,\ldots,S$ (e.g., $S=15$)
    \item \textbf{Regions}: $r = 1,\ldots,R$ (e.g., $R=7$)
    \item \textbf{Obligors}: $i = 1,\ldots,N$
    \item \textbf{Climate drivers}: $k = 1,\ldots,K$ ($K=8$ in our implementation)
\end{itemize}

\textbf{Number of sector--region cells}:
\begin{equation}
M := S \times R \quad \text{(e.g., } M = 105 \text{ for 15 sectors } \times \text{ 7 regions)}
\end{equation}

\textbf{Indexing convention}: Sector-major (regions vary fastest within sector)
\begin{equation}
j = \text{idx}(s,r) := (s-1)R + r
\end{equation}

So $j=1$ is $(s=1, r=1)$, $j=2$ is $(s=1, r=2)$, \ldots, $j=R$ is $(s=1, r=R)$, $j=R+1$ is $(s=2, r=1)$, etc.

\subsection{Factor surface definition}

Define the systematic factor matrix:
\begin{equation}
F \in \mathbb{R}^{S \times R}
\end{equation}

Where $F_{s,r}$ represents the systematic credit condition for sector $s$ in region $r$.

\textbf{Vector form}:
\begin{equation}
\mathbf{f} := \text{vec}(F) \in \mathbb{R}^M
\end{equation}

\textbf{Interpretation}:
\begin{itemize}
    \item $F_{\text{coal,europe}}$ represents systematic credit stress for coal companies in Europe
    \item $F_{\text{agriculture,asia}}$ represents systematic credit stress for agriculture in Asia
    \item These are \textbf{latent random variables} conditional on climate scenario
\end{itemize}

\subsection{Factor decomposition: Scenario mean + residual}

Each cell factor is decomposed as:
\begin{equation}
\mathbf{f} = m(X) + \mathbf{u}
\end{equation}

Where:
\begin{itemize}
    \item $m(X) \in \mathbb{R}^M$: Deterministic scenario-driven mean (climate overlay)
    \item $\mathbf{u} \in \mathbb{R}^M$: Stochastic residual $\sim \mathcal{N}(0, \Sigma_u)$ (non-climate systematic uncertainty)
\end{itemize}

\textbf{Key insight}: Climate fixes the mean; $\Sigma_u$ controls correlated randomness around that mean.

\subsection{Obligor exposure weights (must sum to 1)}

Each obligor $i$ has exposure weights:
\begin{equation}
\beta_i \in \mathbb{R}^M
\end{equation}

With the constraint:
\begin{equation}
\beta_{i,j} \geq 0 \text{ for all } j, \quad \sum_j \beta_{i,j} = 1
\end{equation}

So $\beta_i$ is a \textbf{convex combination} across sector--region cells.

\textbf{Single-cell assignment} (special case):
\begin{itemize}
    \item If obligor $i$ belongs purely to cell $(s(i), r(i))$, then $\beta_i$ has a 1 in position $\text{idx}(s(i),r(i))$ and 0 elsewhere.
\end{itemize}

\textbf{Multi-cell assignment} (diversified exposure):
\begin{itemize}
    \item If obligor $i$ has 60\% revenue from coal in Europe, 40\% from oil in North America:
    \begin{itemize}
        \item $\beta_{\text{idx(coal,europe)}} = 0.6$, $\beta_{\text{idx(oil,NA)}} = 0.4$, all others $= 0$
    \end{itemize}
\end{itemize}

\subsection{Obligor latent variable with asset correlation $\rho$}

We define the obligor's systematic score:
\begin{equation}
S_i := \beta_i^T \cdot \mathbf{f}
\end{equation}

Since different $\beta_i$ imply different systematic variance, we \textbf{standardize}:
\begin{equation}
\text{Var}(S_i) = \beta_i^T \Sigma_u \beta_i
\end{equation}

\begin{equation}
\tilde{S}_i := \frac{S_i}{\sqrt{\beta_i^T \Sigma_u \beta_i}}
\end{equation}

Now $\tilde{S}_i$ has unit variance.

Then we introduce a single global systematic share $\rho \in (0,1)$:
\begin{equation}
Z_i = \sqrt{\rho} \cdot \tilde{S}_i + \sqrt{1-\rho} \cdot \varepsilon_i, \quad \varepsilon_i \sim \mathcal{N}(0,1)
\end{equation}

\textbf{Interpretation of $\rho$}:
\begin{itemize}
    \item $\rho$ is the fraction of variance in $Z_i$ that is systematic for every obligor.
    \item Higher $\rho$ $\Rightarrow$ stronger default clustering and fatter portfolio tails.
    \item Typical range: $\rho \in [0.15, 0.25]$ calibrated to match portfolio tail behavior.
\end{itemize}

\subsection{Default and loss}

Default threshold matches baseline PD:
\begin{equation}
\text{Default}_i = \mathbb{1}\{Z_i < \Phi^{-1}(\text{PD}_i)\}
\end{equation}

Loss:
\begin{equation}
L_i = \text{Default}_i \cdot \text{EAD}_i \cdot \text{LGD}_i
\end{equation}

\begin{equation}
L = \sum_i L_i
\end{equation}

\section{Climate overlay: deterministic scenario drivers $X$}

\subsection{Scenario driver vector $X$}

$X$ is a vector of scenario variables:
\begin{equation}
X \in \mathbb{R}^K \quad (K = 8 \text{ in our implementation})
\end{equation}

It is \textbf{deterministic} for a chosen scenario and horizon. For NGFS ``Net Zero 2050'' year 2050, $X$ is fixed.

\textbf{Deterministic scenario assumption}:
\begin{equation}
X = \mu^{(\text{scenario})}
\end{equation}

No randomness is introduced at the climate-driver level.

\subsection{Climate feature transformation $\phi(X)$}

The raw scenario variables are transformed into standardized features:
\begin{equation}
\phi(X) \in \mathbb{R}^K
\end{equation}

\textbf{Standardization approach}:
\begin{equation}
\phi_k(X) = \frac{X_k - \mu_k^{(\text{baseline})}}{\sigma_k^{(\text{baseline})}}
\end{equation}

Where:
\begin{itemize}
    \item $\mu_k^{(\text{baseline})}$: Mean of driver $k$ under Current Policies scenario
    \item $\sigma_k^{(\text{baseline})}$: Std dev of driver $k$ across scenarios and time horizons
\end{itemize}

This ensures $\phi(X)$ represents ``standard deviations away from business-as-usual''.

\subsection{Cell-level mean systematic factor}

The deterministic scenario-driven mean for each cell is:
\begin{equation}
m(X) = \alpha + W \cdot \phi(X)
\end{equation}

Where:
\begin{itemize}
    \item $\alpha \in \mathbb{R}^M$: Baseline mean (typically zero vector after centering)
    \item $W \in \mathbb{R}^{M \times K}$: Sensitivity matrix (``climate-to-credit weights'')
    \item $W_{j,k}$: How sensitive cell $j$ is to standardized driver $k$
\end{itemize}

\subsection{Hybrid parameterization of $W$ (governance-friendly)}

Instead of calibrating all $M \times K = 840$ weights independently, we use a structured approach:

\begin{equation}
W_{j,k} = \lambda_k \cdot S_{s(j),k} \cdot R_{r(j),k}
\end{equation}

Where:
\begin{itemize}
    \item $\lambda_k \in \mathbb{R}$: Global magnitude parameter for driver $k$ (8 parameters)
    \item $S_{s,k} \in [0,1]$: Sector $s$ sensitivity to driver $k$ (15 $\times$ 8 = 120 parameters)
    \item $R_{r,k} \in [0,1]$: Region $r$ sensitivity to driver $k$ (7 $\times$ 8 = 56 parameters)
\end{itemize}

Total parameters: $8 + 120 + 56 = 184$ (vs 840 independent weights)

This factorization allows us to say:
\begin{itemize}
    \item ``Coal sector is maximally exposed to carbon price'' ($S_{\text{coal,carbon}} = 1.0$)
    \item ``Europe has strong carbon policy exposure'' ($R_{\text{europe,carbon}} = 0.9$)
    \item ``Carbon price transmission is moderately strong'' ($\lambda_{\text{carbon}} = 0.8$)
\end{itemize}

\section{Residual systematic uncertainty: $\Sigma_u$ and correlation structure}

\subsection{The residual covariance matrix}

The stochastic residual captures non-climate systematic uncertainty:
\begin{equation}
\mathbf{u} \sim \mathcal{N}(0, \Sigma_u)
\end{equation}

Where $\Sigma_u \in \mathbb{R}^{M \times M}$ is the residual covariance matrix that encodes:
\begin{itemize}
    \item Inter-sector co-movement
    \item Inter-region co-movement
    \item Portfolio tail thickness (together with $\rho$)
    \item Diversification benefits
\end{itemize}

\textbf{Key design principle}: We \textbf{cannot estimate climate-credit correlations robustly} from historical data. Instead, we:
\begin{itemize}
    \item Keep climate deterministic (no correlation parameters needed)
    \item Put stochastic correlations into $\Sigma_u$
    \item \textbf{Calibrate $\Sigma_u$ from observable market data} (equity sector/region indices)
\end{itemize}

\subsection{Practical construction of $\Sigma_u$ from market data}

Instead of a full free $M \times M$ matrix (5,460 parameters), build it from sector and region correlation blocks:

\textbf{Step 1}: Estimate sector and region correlations from market data
\begin{itemize}
    \item $\text{Corr}_S$: Sector correlation matrix ($S \times S$) estimated from sector index returns (e.g., MSCI sector indices)
    \item $\text{Corr}_R$: Region correlation matrix ($R \times R$) estimated from regional index returns (e.g., MSCI regional indices)
\end{itemize}

\textbf{Step 2}: Convert to covariances with chosen volatility scales
\begin{equation}
\Sigma_S = D_S \cdot \text{Corr}_S \cdot D_S
\end{equation}
\begin{equation}
\Sigma_R = D_R \cdot \text{Corr}_R \cdot D_R
\end{equation}

Where $D_S$ and $D_R$ are diagonal matrices with sector/region volatilities (often start with $D_S = D_R = I$ for simplicity).

\textbf{Step 3}: Embed into cell space using membership matrices

Define \textbf{membership matrices}:
\begin{itemize}
    \item $A \in \mathbb{R}^{M \times S}$: Maps each cell to its sector
    \begin{itemize}
        \item $A_{j,s} = 1$ if cell $j$ belongs to sector $s$, else 0
    \end{itemize}
    \item $B \in \mathbb{R}^{M \times R}$: Maps each cell to its region
    \begin{itemize}
        \item $B_{j,r} = 1$ if cell $j$ belongs to region $r$, else 0
    \end{itemize}
\end{itemize}

\textbf{Step 4}: Mix sector and region components with $\omega$

Define \textbf{sector-vs-region mix parameter $\omega \in [0,1]$}:
\begin{itemize}
    \item $\omega$ controls whether residual co-movement is primarily sector-cycle or region-cycle driven
    \item High $\omega$ ($\to 1$): Sector cycles dominate (e.g., energy sector co-moves globally)
    \item Low $\omega$ ($\to 0$): Regional cycles dominate (e.g., geographic macro conditions matter more)
    \item Typical starting point: $\omega = 0.5$ (equal weight)
\end{itemize}

Construct the structured covariance component:
\begin{equation}
\Sigma_{\text{structured}} = \omega \cdot (A \Sigma_S A^T) + (1-\omega) \cdot (B \Sigma_R B^T)
\end{equation}

\textbf{Step 5}: Mix structured vs cell-specific with $\eta$

Define \textbf{structured-vs-cell share parameter $\eta \in [0,1]$}:
\begin{itemize}
    \item $\eta$ controls how much of residual variance is structured (via $\Sigma_{\text{structured}}$) vs cell-specific
    \item High $\eta$ ($\to 1$): Strong systematic clustering, less diversification benefit
    \item Low $\eta$ ($\to 0$): More cell-specific noise, greater diversification benefit
    \item Typical starting point: $\eta = 0.7$ (moderately systemic)
\end{itemize}

Construct final residual covariance:
\begin{equation}
\Sigma_u = \eta \cdot \Sigma_{\text{structured}} + (1-\eta) \cdot I_M
\end{equation}

Expanding:
\begin{equation}
\boxed{\Sigma_u = \eta \cdot [\omega \cdot (A \Sigma_S A^T) + (1-\omega) \cdot (B \Sigma_R B^T)] + (1-\eta) \cdot I_M}
\end{equation}

\subsection{Calibration knobs}

\begin{itemize}
    \item $\text{Corr}_S$, $\text{Corr}_R$: Estimated from data (MSCI sector/regional indices, 5-year rolling correlations)
    \item $D_S$, $D_R$: Volatility scales (typically set to $I$ for unit volatility)
    \item $\omega$: Sector-vs-region mix ($0$ = all region, $1$ = all sector, $0.5$ = equal)
    \item $\eta$: Structured-vs-cell share ($0$ = all cell-specific, $1$ = all structured)
\end{itemize}

\subsection{Induced correlation structure}

The construction induces:

\textbf{Within-sector correlation} (same sector, different regions):
\begin{equation}
\text{Corr}(u_{s,r}, u_{s,r'}) = \frac{\eta \cdot \omega \cdot \text{Sector variance component}}{\text{Total variance}}
\end{equation}

Controlled by $\eta$ (strength of structure) and $\omega$ (sector vs region weight).

\textbf{Within-region correlation} (different sectors, same region):
\begin{equation}
\text{Corr}(u_{s,r}, u_{s',r}) = \frac{\eta \cdot (1-\omega) \cdot \text{Region variance component}}{\text{Total variance}}
\end{equation}

Controlled by $\eta$ and $(1-\omega)$.

\textbf{Diagonal variance} (variance of each cell):
\begin{equation}
\text{Var}(u_{s,r}) = \eta \cdot [\omega \cdot (A \Sigma_S A^T)_{s,r} + (1-\omega) \cdot (B \Sigma_R B^T)_{s,r}] + (1-\eta)
\end{equation}

With $D_S = D_R = I$, the diagonal is approximately $\eta + (1-\eta) = 1$ (unit variance cells).

\section{The calibration knobs: what you must choose/estimate}

\subsection{8 key parameters}

The model has 8 main calibration ``knobs'':

\begin{enumerate}
    \item \textbf{S matrix} ($15 \times 8 = 120$ parameters): Sector sensitivity scores $S_{s,k} \in [0,1]$
    \item \textbf{R matrix} ($7 \times 8 = 56$ parameters): Region sensitivity scores $R_{r,k} \in [0,1]$
    \item \textbf{$\lambda$ vector} (8 parameters): Global magnitude parameters $\lambda_k$
    \item \textbf{$\text{Corr}_S$} ($15 \times 15$ matrix): Sector correlation (estimated from MSCI sector indices)
    \item \textbf{$\text{Corr}_R$} ($7 \times 7$ matrix): Region correlation (estimated from MSCI regional indices)
    \item \textbf{$\omega$} (1 parameter): Sector-vs-region mix $\in [0,1]$
    \item \textbf{$\eta$} (1 parameter): Structured-vs-cell share $\in [0,1]$
    \item \textbf{$\rho$} (1 parameter): Asset correlation $\in [0,1]$
\end{enumerate}

\section{Simulation algorithm (single horizon)}

Given scenario $X = \mu^{(\text{scenario})}$:

\begin{enumerate}
    \item Compute features $\phi(X)$
    \item Compute deterministic mean $m(X) = \alpha + W \cdot \phi(X)$
    \item Draw residual systematic factor $\mathbf{u} \sim \mathcal{N}(0, \Sigma_u)$
    \item Form $\mathbf{f} = m(X) + \mathbf{u}$
    \item For each obligor $i$:
    \begin{enumerate}
        \item Compute systematic score $S_i = \beta_i^T \mathbf{f}$
        \item Standardize $\tilde{S}_i = S_i / \sqrt{\beta_i^T \Sigma_u \beta_i}$
        \item Draw $\varepsilon_i \sim \mathcal{N}(0,1)$
        \item Compute $Z_i = \sqrt{\rho} \tilde{S}_i + \sqrt{1-\rho} \varepsilon_i$
        \item Set $\text{Default}_i = \mathbb{1}\{Z_i < \Phi^{-1}(\text{PD}_i)\}$
        \item Compute $L_i = \text{Default}_i \cdot \text{EAD}_i \cdot \text{LGD}_i$
    \end{enumerate}
    \item Aggregate $L = \sum_i L_i$
\end{enumerate}

Repeat steps (3)--(6) to obtain the loss distribution under the scenario.

\section{Notation Reference}

\begin{table}[h]
\centering
\begin{tabular}{ll}
\toprule
\textbf{Symbol} & \textbf{Description} \\
\midrule
$X$ & Deterministic climate/scenario variable vector \\
$\phi(X)$ & Standardized climate features \\
$F, \mathbf{f}$ & Sector--region factor matrix (and vectorized form) \\
$m(X)$ & Scenario-driven mean systematic factor \\
$\mathbf{u}$ & Stochastic residual systematic factor \\
$\Sigma_u$ & Residual covariance matrix \\
$W$ & Climate-to-credit sensitivity matrix ($M \times K$) \\
$S$ & Sector sensitivity matrix ($15 \times 8$) \\
$R$ & Region sensitivity matrix ($7 \times 8$) \\
$\lambda$ & Global magnitude vector (8 scalars) \\
$\text{Corr}_S$ & Sector correlation matrix ($15 \times 15$) \\
$\text{Corr}_R$ & Region correlation matrix ($7 \times 7$) \\
$\omega$ & Sector-vs-region mix parameter $\in [0,1]$ \\
$\eta$ & Structured-vs-cell share parameter $\in [0,1]$ \\
$\rho$ & Asset correlation parameter $\in [0,1]$ \\
$\beta_i$ & Obligor $i$ exposure weights (convex combination over cells) \\
$Z_i$ & Obligor $i$ latent variable \\
$\text{PD}_i$ & Obligor $i$ probability of default \\
$\text{EAD}_i$ & Obligor $i$ exposure at default \\
$\text{LGD}_i$ & Obligor $i$ loss given default \\
$L$ & Portfolio loss \\
\bottomrule
\end{tabular}
\end{table}

\end{document}
